% gerado por ferramentas/md2tex_enunciado.py a partir de questoes/q130.md (não editar à mão)
Um jardineiro está irrigando as plantas de um jardim com uma mangueira. Ele estica completamente a mangueira para se aproximar de algumas flores mais afastadas, sem, no entanto, alcançá-las. Daí, ocorre-lhe que obstruir a extremidade da mangueira com o polegar faz com que o jato de água atinja uma distância maior, comparativamente com a situação sem essa ação, permitindo-lhe molhar aquelas flores mais afastadas.\par
Essa alteração no alcance do jato de água deve-se ao aumento do(a)\par
\begin{alternativas}
  \item velocidade da água lançada.
  \item densidade da água lançada.
  \item volume da água lançada.
  \item massa da água lançada.
  \item vazão da água lançada.
\end{alternativas}
